%-----------------------------------------------------------------
% BAB I PENDAHULUAN
%
% Struktur sesuai panduan penulisan tesis:
%   1.1 Latar Belakang
%   1.2 Rumusan Masalah
%   1.3 Tujuan Penelitian
%   1.4 Manfaat Penelitian
%-----------------------------------------------------------------
\chapter{PENDAHULUAN}

% Struktur bab mengikuti panduan penulisan tesis.

\section{Latar Belakang}

Tuliskan latar belakang penelitian beserta data-data empiris yang
melatarbelakangi. Tutup dengan celah penelitian (research gap) yang akan
diisi oleh penelitian ini, serta rujukan dari `daftar-pustaka.bib`
dengan perintah \cite{buku2020}.

\section{Rumusan Masalah}

Tuliskan rumusan masalah penelitian.

\section{Tujuan Penelitian}

Tuliskan tujuan penelitian.

\section{Manfaat Penelitian}

Tuliskan manfaat penelitian bagi ilmu pengetahuan, teknologi, masyarakat,
dan komunitas ilmiah.